\documentclass{beamer}
\usepackage[utf8]{inputenc}

\usetheme{Berkeley}
\usecolortheme{default}

\setbeamertemplate{footline}[frame number]
\beamertemplatenavigationsymbolsempty

\usepackage[most]{tcolorbox}

\title{Project Workflows and Writing Tools}
\subtitle{From Ideas to Submission-Ready Manuscripts}
\author[]{IB 514: Scientific Writing}
\date[]{}

\AtBeginSection[]
{
  \begin{frame}
    \frametitle{Overview}
    \tableofcontents[currentsection]
  \end{frame}
}



\begin{document}

%------------------------------------------------
\AtBeginSection[]
{
  \begin{frame}
    \frametitle{Overview}
    \tableofcontents[currentsection]
  \end{frame}
}

%------------------------------------------------
\begin{frame}
  \titlepage
\end{frame}

% -------------------------------------------------


\begin{frame}{Framing the Session: Why Workflows Matter}

\small

\textbf{Central Premise:}
    High-quality writing emerges from high-quality systems, not last-minute effort.


\vspace{0.5em}

\textbf{Three Guiding Questions}
\begin{enumerate}
    \item How can I reduce friction between ideas and written text?
    \item How can I prevent structural chaos as projects grow?
    \item How can I ensure my workflow supports collaboration and revision?
\end{enumerate}

\vspace{0.5em}

\textbf{Session Objectives}
\begin{itemize}
    \item Design a scalable project structure
    \item Identify points where inefficiency may emerge
    \item Evaluate writing tools strategically
\end{itemize}

\vspace{0.5em}

\begin{block}{\textbf{Key Insight:} Workflow design is a methodological decision }
    It shapes clarity, reproducibility, and speed.
\end{block}

\end{frame}

% -------------------------------------------------

\section{Project Structure and Organization}

\begin{frame}{Project Structure}

\vspace{0.5em}

\textbf{Recommended Folder Architecture}

    \begin{itemize}
        \item \texttt{/notes} – reading notes, idea sketches, meeting notes
        \item \texttt{/data}
            \begin{itemize}
                \item \texttt{/raw} (read-only)
                \item \texttt{/processed}
            \end{itemize}
        \item \texttt{/code} – scripts, functions, notebooks
        \item \texttt{/results} – exploratory outputs, intermediary results
        \item \texttt{/tables} - publication-ready tables\footnote{Exported if using \LaTeX}
        \item \texttt{/figures} – exported publication-ready figures
        \item \texttt{/manuscript} – bibliography, drafts and versions\footnote{If not using version control software}
    \end{itemize}

\end{frame}

% ---------------------------------------------

\begin{frame}{Project Structure}

\textbf{Best Practices}
\begin{itemize}
    \item Explanatory README file in root directory
    \item Separate raw and cleaned/processed data
    \item Separate exploratory from finalized results (figs and tables)
    \item Consistent naming conventions
\end{itemize}

\vspace{0.5em}

\textbf{File Naming Conventions}

\begin{itemize}
    \item Consistent ordering:
    \begin{itemize}
        \item \texttt{survey\_raw\_2025.csv}
        \item \texttt{survey\_cleaned\_2025.csv}
    \end{itemize}
    \textbf{Consistent capitalization:}
    \begin{itemize}
        \item \texttt{Upperstart} vs. \texttt{lowercase}
        \item \texttt{CamelCase} vs. \texttt{start\_end}
    \end{itemize}
    \item Use descriptive, machine-readable names:
    \begin{itemize}
        \item Avoid: \texttt{final\_real\_final2.docx}
        \item Prefer: \texttt{results/regression\_modelA.csv}
    \end{itemize}
\end{itemize}

\end{frame}


%------------------------------------------------
\begin{frame}{Version Control}

\textbf{Why Version Control?}
\begin{itemize}
    \item Prevents loss of work
    \item Eases identification and fixing of problems/bugs
    \item Avoid concern over losing removed text
    \item Enables collaboration
    \item Documents intellectual development
\end{itemize}


\textbf{Approaches}
\begin{itemize}
    \item Manual versioning 
        \begin{itemize}
            \item Use ISO dates: \texttt{2026-02-15\_analysis.R}
            \item Version numbers: \texttt{manuscript\_v03.tex}
            \item Coauthor comments: \texttt{manuscript\_v03\_MN.tex}
        \end{itemize}
    \item Cloud version history (synced folders, e.g., Dropbox, Box)
    \item Git-based workflows (e.g., Git with GitHub)
\end{itemize}


\end{frame}

%------------------------------------------------
\section{Writing Workflows}

\begin{frame}{Writing Workflow: From Ideas to Manuscript}
\textbf{Linear model (non-existent):}
\[
\text{Ideas} \rightarrow \text{Notes} \rightarrow \text{Data} \rightarrow \text{Code} \rightarrow \text{Figures} \rightarrow \text{Manuscript}
\]

\vspace{0.5em}
\begin{block}
    {\textbf{Key Insight:} The workflow is iterative, not strictly linear.}
    Most time is spent revising.
\end{block}
\vspace{0.5em}

\textbf{Common Transition Points}
\begin{itemize}
    \item Notes $\rightarrow$ structured research questions
    \item Code $\rightarrow$ validated results
    \item Clarity $\rightarrow$ candidate figures
    \item Figures $\rightarrow$ narrative refinement
\end{itemize}

\begin{block}{\textbf{Reminder:}}
    Draft sections early (especially Methods and Results).
\end{block}

\end{frame}

%------------------------------------------------

\begin{frame}{Pre-Drafting Structural Design}

\begin{block}{\textbf{Core Principle:}}
    Design structure deliberately; diagnose structure explicitly.
\end{block}

\textbf{Storyboarding}
\begin{itemize}
    \item Construct the argument visually before drafting prose.
    \item Create one slide, page, or card per figure or major claim.
    \item For each, specify:
    \begin{itemize}
        \item The core claim (1 sentence)
        \item The evidence shown
        \item Why it matters for the overall argument
    \end{itemize}
    \item Arrange in logical order to test narrative flow.
    \item Identify missing transitions or unsupported claims early.
\end{itemize}

\textbf{Goal:}
Ensure overall logical coherence before drafting prose\footnote{But remain open to change!}

\end{frame}

% --------------------------------------------
\begin{frame}{Post-Drafting Structural Diagnosis}

\begin{block}{\textbf{Core Principle:}}
    Design structure deliberately; diagnose structure explicitly.
\end{block}

\textbf{Reverse Outlining}
\begin{itemize}
    \item Extract 1–2 bullet-point claims per paragraph.
    \item Ignore prose quality; focus only on argumentative function.
    \item Ask:
    \begin{itemize}
        \item Does each paragraph have a clear purpose?
        \item Does the sequence of bullets form a coherent argument?
        \item Are there redundancies or logical gaps?
    \end{itemize}
    \item Revise structure before revising sentences.
\end{itemize}

\textbf{Goal:}
Ensure both internal and overall logical coherence after drafting prose.

\end{frame}




%------------------------------------------------
\section{Journal Article Types}

%------------------------------------------------
\begin{frame}{IMRaD (Introduction–Methods–Results–Discussion)}
    \fontsize{9pt}{10pt}\selectfont

\begin{enumerate}
    \item \textbf{Methods} first (most structurally constrained).
    \item \textbf{Results} second (figure-driven; claim-evidence alignment).
    \item \textbf{Introduction} after results clarify true contribution.
    \item \textbf{Discussion} last to interpret, generalize, and situate.
\end{enumerate}

\vspace{0.4em}

\textbf{Introduction (Why?)}
\begin{itemize}
    \item Move from broad problem to specific knowledge gap.
    \item Critically evaluate literature (don't summarize passively).
    \item End with explicit research questions, hypotheses, or key finding.
\end{itemize}

\textbf{Methods (How?)}
\begin{itemize}
    \item Provide enough detail for replication.
    \item Use clear subheadings aligned with research questions.
    \item Justify key analytical decisions.
\end{itemize}

\textbf{Results (What?)}
\begin{itemize}
    \item Organize around figures/tables.
    \item Lead each subsection with the main finding, align with Methods sections
    \item Avoid interpretation beyond minimal context.
\end{itemize}

\textbf{Discussion (So what?)}
\begin{itemize}
    \item Interpret findings relative to hypotheses.
    \item Compare with prior literature.
    \item Address limitations and articulate broader implications.
\end{itemize}

\end{frame}

%------------------------------------------------
\begin{frame}{Methods-at-the-End Papers}
\fontsize{9pt}{10pt}\selectfont


\textbf{Design a Figure-Driven Narrative}
\begin{itemize}
    \item Determine the logical sequence of core results.
    \item Ensure each figure advances the central claim.
    \item Remove peripheral analyses from main text.
\end{itemize}

\textbf{Front-Load Essential Context}
\begin{itemize}
    \item Provide minimal methodological orientation in Results.
    \item Define key variables and datasets before presenting outcomes.
\end{itemize}

\textbf{Write Results as a Coherent Argument}
\begin{itemize}
    \item Each subsection should answer a focused question.
    \item Maintain conceptual momentum between figures.
\end{itemize}

\textbf{Develop Detailed Methods Section at End}
\begin{itemize}
    \item Provide full reproducibility detail.
    \item Use clear subsections for sampling, measurements, analyses.
    \item Anticipate reviewer questions and address proactively.
\end{itemize}

\textbf{Key:}
Requires exceptionally clear storytelling

\end{frame}

% -----------------------------------------
\begin{frame}{Literature Reviews}
\fontsize{9pt}{10pt}\selectfont
    
\begin{itemize}
    \item \textbf{Define scope precisely}
    \begin{itemize}
        \item What question is the review answering?
        \item What time period, disciplines, or methodologies are included?
    \end{itemize}
    \item \textbf{Develop a transparent search strategy}
    \begin{itemize}
        \item Databases used, search terms, inclusion/exclusion criteria.
        \item Consider a flow diagram for study selection.
    \end{itemize}
    \item \textbf{Organize by conceptual structure (not chronology)}
    \begin{itemize}
        \item Thematic clusters
        \item Competing hypotheses
        \item Methodological approaches
    \end{itemize}
    \item \textbf{Synthesize within sections}
    \begin{itemize}
        \item Compare studies directly.
        \item Highlight agreements, disagreements, and methodological drivers.
    \end{itemize}
    \item \textbf{Conclude with implications}
    \begin{itemize}
        \item Identify unresolved questions.
        \item Clarify future research directions.
    \end{itemize}
\end{itemize}

\end{frame}


\begin{frame}{Synthesis}
\fontsize{9pt}{10pt}\selectfont

\begin{itemize}
    \item \textbf{Begin with a conceptual problem }
    \begin{itemize}
        \item What tension, paradox, or fragmentation exists?
    \end{itemize}
    \item \textbf{Abstract upward }
    \begin{itemize}
        \item Move from specific findings to higher-level principles.
    \end{itemize}
    \item \textbf{Integrate across domains }
    \begin{itemize}
        \item Combine theories, datasets, or methods not typically linked.
    \end{itemize}
    \item \textbf{Develop a unifying framework }
    \begin{itemize}
        \item Conceptual diagram or model.
        \item Clear definitions of key constructs.
    \end{itemize}
    \item \textbf{Demonstrate explanatory gain }
    \begin{itemize}
        \item Show how the synthesis resolves contradictions or generates predictions.
    \end{itemize}
\end{itemize}

\end{frame}

\begin{frame}{Data Papers}
\begin{itemize}
    \item Clarify dataset purpose and potential reuse scenarios.
    \item Provide full documentation:
    \begin{itemize}
        \item Sampling design, instrumentation, preprocessing steps.
        \item File structure and variable-level metadata.
    \end{itemize}
    \item Demonstrate quality control
    \begin{itemize}
        \item Validation checks, error screening, uncertainty description.
    \end{itemize}
    \item Include example use cases
    \begin{itemize}
        \item Illustrative analyses showing dataset functionality.
    \end{itemize}
\end{itemize}

\end{frame}


\section{Writing Tools}
\begin{frame}[allowframebreaks]{Writing Tools}
    \fontsize{9pt}{8pt}\selectfont
    \begin{columns}[t]
        \begin{column}{0.48\textwidth}
            \begin{block}{Word / Pages / Google Docs}
                \textbf{Pros:}
                \begin{itemize}
                    \item Widely used; track-changes
                    \item Real-time cloud collaboration
                \end{itemize}
                \textbf{Cons:}
                \begin{itemize}
                    \item No or poor version control
                    \item Formatting instability
                \end{itemize}
            \end{block}
            
            \begin{block}{\LaTeX}
                \textbf{Pros:}
                \begin{itemize}
                    \item Lightweight, plain text
                    \item Superior math, refs, tables
                    \item Auto-numbering \& insertion
                \end{itemize}
                \textbf{Cons:}
                \begin{itemize}
                    \item Learning curve
                    \item No real-time WYSIWYG
                    \item Harder for collaborators
                    \item No built-in track-changes
                \end{itemize}
            \end{block}
        \end{column}

        \begin{column}{0.48\textwidth}
            \begin{block}{Markdown}
                \textbf{Pros:}
                \begin{itemize}
                    \item Lightweight, plain text
                    \item Future-proof format
                    \item Easy code integration
                \end{itemize}
                \textbf{Cons:}
                \begin{itemize}
                    \item Toolchain setup required
                    \item Limited image handling
                \end{itemize}
            \end{block}

            \begin{block}{Scrivener}
                \textbf{Pros:}
                \begin{itemize}
                    \item Research \& reference binder
                    \item Snapshots for safe revisions
                    \item Great for nonlinear drafting
                \end{itemize}
                \textbf{Cons:}
                \begin{itemize}
                    \item Not for collaborative work
                    \item Proprietary file format
                    \item Export cleanup required
                \end{itemize}
            \end{block}
        \end{column}
    \end{columns}
\end{frame}


%------------------------------------------------
%------------------------------------------------
%------------------------------------------------
%------------------------------------------------
\begin{frame}{Additional Tools}

\fontsize{10pt}{10pt}\selectfont

\textbf{Dictation Tools}

e.g, MS Speech Recognition,  macOS Dictation, Otter.ai {\tiny (300 min./mo.)}

\begin{columns}[T, totalwidth=\textwidth]
    \column{0.48\textwidth}
    {\fontsize{8pt}{9pt}\selectfont
    \textbf{Pros:}
    \begin{itemize}
        \item Rapidly capture ideas, reduce typing fatigue
        \item Lowers activation energy to overcome writing inertia
    \end{itemize}
    }

    \column{0.48\textwidth}
    {\fontsize{8pt}{9pt}\selectfont
    \textbf{Cons:}
    \begin{itemize}
        \item Substantial editing (typos, grammar, clarity, formatting)
        \item Misinterpretation of technical or domain-specific terms
    \end{itemize}
    }
\end{columns}

\vspace{0.5em}
\noindent\rule{\textwidth}{1pt}
\vspace{0.5em}

\textbf{Large Language Models (LLMs)}

    e.g., ChatGPT, Claude, Scite, Elicit

\begin{columns}[T, totalwidth=\textwidth]
    \column{0.48\textwidth}
    {\fontsize{8pt}{9pt}\selectfont
    \textbf{Pros:}
    \begin{itemize}
        \item Assist in drafting sentences, paragraphs, and structured outlines
        \item Suggest structural improvements, phrasing, and organization
        \item Literature summaries or conceptual scaffolds
    \end{itemize}
    }

    \column{0.48\textwidth}
    {\fontsize{8pt}{9pt}\selectfont
    \textbf{Cons:}
    \begin{itemize}
        \item Risk of (over)reliance, inhibiting original thought
        \item Factual inaccuracies, limited domain knowledge or literature access
        \item Generic or formulaic writing
        \item Environmental impact
    \end{itemize}
    }
\end{columns}

\end{frame}
%------------------------------------------------
%------------------------------------------------
%------------------------------------------------

%------------------------------------------------
\section{Figures}

\begin{frame}{Conceptual and Data Figures}
\textbf{Conceptual Figures}
\begin{itemize}
    \item Clarify theoretical frameworks
    \item Summarize complex relationships
\end{itemize}

\textbf{Data Figures}
\begin{itemize}
    \item Present core results
    \item Should be interpretable independently
\end{itemize}

\textbf{Manual vs. Code-Based Creation}
\begin{itemize}
    \item Manual (e.g., Illustrator): high aesthetic control
    \item Code-based: reproducible and scalable
\end{itemize}

\textbf{File Formats Commonly Requested}
\begin{itemize}
    \item PDF – typically sufficient
    \item TIFF – high-resolution images
    \item EPS – vector graphics
    \item Figure resolution (e.g., 300–600 dpi)
\end{itemize}

\end{frame}

%------------------------------------------------
\section{Final thoughts}
\begin{frame}{Workflow Integration}
\textbf{Align structure, tools, and formatting strategy.}

\begin{itemize}
    \item Organize for clarity early to reduce friction later.
    \item Choose tools that match writing needs.
    \item Separate development and writing from formatting.
\end{itemize}

\begin{block}{\textbf{Goal:}}
    Efficient, transparent, and reproducible writing workflow.
\end{block}
\end{frame}

%------------------------------------------------
\end{document}